\section{Installation und Betrieb}
\label{sec:installation}

SignalReport ist als einzelne ausführbare JAR-Datei konzipiert und benötigt lediglich
eine Java-21-Laufzeitumgebung. Dieses Kapitel beschreibt die verschiedenen
Installationswege und den Betrieb als Hintergrund-Dienst.

\subsection{Voraussetzungen}

\begin{itemize}
  \item \textbf{Java 21+} -- OpenJDK (Temurin) oder Oracle JDK
  \item \textbf{Betriebssystem:} Windows 10/11, macOS 13+, Linux (systemd), Synology DSM 7+
  \item \textbf{Arbeitsspeicher:} mind.\ 256\,MB frei (empfohlen: 512\,MB)
  \item \textbf{Speicherplatz:} ca.\ 50\,MB für die Anwendung; Datenbank wächst mit der
    Laufzeit (ca.\ 1\,MB pro Tag bei 10-Sekunden-Intervall)
  \item \textbf{Netzwerk:} Port 4567 (TCP) muss frei sein
\end{itemize}

\subsection{Schnellstart (portabel)}

Der einfachste Weg, SignalReport zu starten:

\begin{lstlisting}[language=bash, caption={Portabler Start}]
java -jar signalreport.jar
\end{lstlisting}

Anschließend ist das Web-Interface unter \texttt{http://localhost:4567} erreichbar.
Beim ersten Start wird automatisch der Setup-Assistent angezeigt.

\textbf{Hinweis:} Bei dieser Variante läuft SignalReport nur, solange das Terminal-Fenster
geöffnet bleibt. Für den Dauerbetrieb empfiehlt sich die Installation als Dienst.

\subsection{Installation als Windows-Dienst}
\label{subsec:install-windows}

Das Projekt enthält ein Installations-Skript unter
\texttt{deployment/windows/install.bat}, das SignalReport mittels Apache Commons Daemon
(procrun) als Windows-Dienst einrichtet.

\subsubsection{Ablauf}

\begin{enumerate}
  \item Skript als Administrator ausführen (Rechtsklick $\rightarrow$ ``Als Administrator
    ausführen'')
  \item Automatische Prüfung: Administrator-Rechte und Java-Installation
  \item Erstellung der Verzeichnisse:
    \begin{itemize}
      \item \texttt{\%ProgramFiles\%\textbackslash SignalReport} -- Programmdateien
      \item \texttt{\%ProgramData\%\textbackslash SignalReport} -- Daten und Logs
    \end{itemize}
  \item Download von Apache Commons Daemon 1.5.1 (\texttt{prunsrv.exe})
  \item Registrierung des Windows-Dienstes ``SignalReport'' (Autostart, LocalSystem)
  \item Erstellung einer Desktop-Verknüpfung
  \item Einrichtung einer Firewall-Regel für Port 4567
\end{enumerate}

\subsubsection{Verwaltung}

\begin{lstlisting}[language=bash, caption={Windows-Dienst verwalten}]
# Status pruefen
sc query SignalReport

# Dienst stoppen/starten
net stop SignalReport
net start SignalReport

# Service-Manager oeffnen (GUI)
"%ProgramFiles%\SignalReport\prunmgr.exe"
\end{lstlisting}

Die Deinstallation erfolgt über \texttt{deployment/windows/uninstall.bat} (ebenfalls als
Administrator).

\subsection{Installation unter Linux (systemd)}
\label{subsec:install-linux}

Das Skript \texttt{deployment/macos-linux/install.sh} erkennt automatisch die Plattform
und erstellt unter Linux einen systemd-Service.

\subsubsection{Ablauf}

\begin{enumerate}
  \item Ausführung mit \texttt{sudo ./install.sh}
  \item Erstellung eines System-Benutzers \texttt{signalreport} (kein Login)
  \item Verzeichnisstruktur:
    \begin{itemize}
      \item \texttt{/opt/signalreport/} -- Programmdateien
      \item \texttt{/var/lib/signalreport/} -- Datenbank und Konfiguration
      \item \texttt{/var/log/signalreport/} -- Log-Dateien
    \end{itemize}
  \item Optional: Kompilierung von jsvc (Apache Commons Daemon) für professionellen
    Daemon-Modus; Fallback auf direkten Java-Start
  \item Erstellung und Aktivierung der systemd-Unit \texttt{signalreport.service}
\end{enumerate}

\subsubsection{Verwaltung}

\begin{lstlisting}[language=bash, caption={systemd-Dienst verwalten}]
# Status pruefen
sudo systemctl status signalreport

# Dienst stoppen/starten/neustarten
sudo systemctl stop signalreport
sudo systemctl start signalreport
sudo systemctl restart signalreport

# Logs anzeigen
sudo journalctl -u signalreport -f
\end{lstlisting}

\subsection{Installation unter macOS (launchd)}
\label{subsec:install-macos}

Das gleiche Skript (\texttt{install.sh}) erkennt macOS und erstellt einen launchd-Service.

\subsubsection{Unterschiede zu Linux}

\begin{itemize}
  \item Installationsverzeichnis: \texttt{/Applications/SignalReport/}
  \item Datenverzeichnis: \texttt{\textasciitilde/Library/Application Support/SignalReport/}
  \item Service-Datei: \texttt{/Library/LaunchDaemons/com.signalreport.service.plist}
  \item Desktop-Verknüpfung: \texttt{.webloc}-Datei auf dem Schreibtisch
  \item Läuft als angemeldeter Benutzer (nicht als separater Systembenutzer)
\end{itemize}

\subsubsection{Verwaltung}

\begin{lstlisting}[language=bash, caption={launchd-Dienst verwalten}]
# Dienst stoppen
sudo launchctl stop com.signalreport.service

# Dienst starten
sudo launchctl start com.signalreport.service

# Dienst komplett entladen/laden
sudo launchctl unload /Library/LaunchDaemons/com.signalreport.service.plist
sudo launchctl load /Library/LaunchDaemons/com.signalreport.service.plist
\end{lstlisting}

\subsection{Synology NAS (Docker)}

Auf einem Synology NAS mit DSM 7+ und installiertem Container Manager kann SignalReport
als Docker-Container betrieben werden. Ein Dockerfile ist im Verzeichnis
\texttt{deployment/docker/} vorbereitet.

\subsection{Konfiguration}
\label{subsec:konfiguration}

Nach dem ersten Start wird automatisch eine \texttt{config.json} im Arbeitsverzeichnis
erstellt. Alle Einstellungen können über das Web-Interface geändert werden; alternativ
kann die Datei manuell bearbeitet werden (bei gestopptem Dienst).

\begin{lstlisting}[language=json, caption={Beispiel config.json (Auszug)}]
{
  "measurement": {
    "intervalSeconds": 10,
    "targets": {
      "ping": "8.8.8.8",
      "dns": "google.com",
      "http": "https://example.com"
    }
  },
  "maintenanceWindow": {
    "enabled": true,
    "startHour": 4, "startMinute": 0,
    "endHour": 4, "endMinute": 10
  },
  "userInfo": {
    "provider": "Magenta",
    "customerId": "08154711",
    "userName": "Max Mustermann"
  },
  "theme": {
    "darkMode": false
  }
}
\end{lstlisting}

\subsection{Backup und Migration}

Für ein vollständiges Backup genügt es, zwei Dateien/Verzeichnisse zu sichern:

\begin{enumerate}
  \item \texttt{config.json} -- Alle Einstellungen (Ziele, Passwort-Hashes, Theme)
  \item \texttt{data/} -- H2-Datenbank mit allen Messdaten
\end{enumerate}

Zur Migration auf ein anderes System: SignalReport installieren, diese beiden
Dateien/Verzeichnisse kopieren und den Dienst starten. Die Datenbank und Konfiguration
werden automatisch erkannt.
